\section{Comprobación exacta de las construcciones finitas}
\label{sec:exact-reproduction}

La reproducción distingue la demostración general y la evaluación de
una especialización. Las proposiciones precedentes utilizan identidades
algebraicas, positividad e hipótesis de convergencia explícitas. Sus
controles computacionales se realizan con enteros y fracciones racionales;
el desarrollo reticular conserva, además, su representación exacta en
la extensión algebraica declarada. Los datos del registro son salidas
expresadas por la construcción anterior. Las ejecuciones de esta
sección comprueban su transporte geométrico y operatorio, y conservan
por separado la procedencia de su generación.

El primer control recupera las separaciones a partir de los menores,
comprueba la relación de Plücker en las triangulaciones y compara
cambios de coordenatización con deformaciones efectivas. Un cambio invertible
de las filas multiplica todos los menores por el mismo determinante;
los cocientes marcados permanecen iguales. Un reescalado independiente
de las columnas conserva ciertas razones dobles y altera, en general,
las separaciones reconstruidas. Evaluar ambas operaciones sobre los
mismos datos proporciona un control negativo de la equivalencia
que se afirma.

El segundo control subdivide las longitudes por nueve, conserva las
marcas heredadas y verifica las identidades de promedios y de Gram.
Los determinantes se calculan por eliminación exacta y se contrastan
con sus productos de Vandermonde. La positividad en cualquier
dimensión finita procede de la fórmula demostrada; los casos finitos
ejecutados comprueban la implementación de esa fórmula, sus índices
y sus normalizaciones. Las medidas de conteo y las medidas
cronológicas se mantienen separadas para que una igualdad de
momentos no dependa de cambiar silenciosamente la normalización.

El tercer control reproduce los ocho primeros momentos mediante
el operador de Jacobi de orden cuatro. La matriz simétrica se
representa de forma racional en la base de polinomios mónicos;
su métrica es \(\operatorname{diag}(h_0,\ldots,h_3)\). Se
comprueba la autoadjunción en esa métrica, la recurrencia de tres
términos y el acuerdo entre el resolvente matricial y su fracción
continua en \(z=2\). El resultado mantiene explícitos el vector
constante observado y la carga \(Q\), que intervienen en el
teorema de fidelidad espectral marcada.

El cuarto grupo verifica los cambios simplécticos y sus
Hamiltonianos. La derivada respecto del momento recupera la
velocidad; la transformación de Legendre se comprueba como
identidad polinómica o de Laurent. La inercia de una forma
cuadrática y la positividad de un Hamiltoniano cuántico
pertenecen a sus espacios respectivos. La representación
de Clifford se comprueba mediante sus anticonmutadores,
y la composición con el operador de masa se prueba por
el cuadrado del operador de Dirac. Estas igualdades
identifican cada término y sus productos cruzados.

Finalmente se conservan las comprobaciones anteriores de
triangulación, forma canónica, reducción de Schur, congruencia
de masas y construcción reticular. En Yang--Mills, las
proposiciones sobre medida común, refinamiento y reconstrucción
se exponen en el cuerpo con sus premisas y falsadores. Un
control finito de matrices no sustituye el paso al límite
ni la reconstrucción de campos: sus dominios se registran
separadamente en la entrega reproducible.

El paquete de fuentes enlaza cada programa con el enunciado
que evalúa y conserva su salida íntegra. Las condiciones se
comprueban mediante bifurcaciones explícitas, de modo que
permanecen activas al ejecutar el intérprete en modo
optimizado. Se reproducen asimismo en modo aislado, sin
cargar paquetes de usuario. La correspondencia entre
construcción, prueba y ejecución puede inspeccionarse
paso a paso sin convertir la salida de un programa en
una afirmación de alcance superior a la que comprueba.
